\section*{Bab \ref{ch:b3:molecular-evolution} --- Evolusi Molekuler dan Filogenomika}

\begin{solution}{exo:b3:molecular-evolution:1}
$k = \mu$: \omterm{thm:b3:molecular-evolution:neutral}{laju substitusi} netral tiap tapak tiap angkatan sama dengan
laju mutasinya. Tiap angkatan $2N\mu$ mutasi netral baru timbul
dan masing-masing berpeluang $1/2N$ terpancang; sebuah populasi yang lebih besar membuat
lebih banyak mutasi lalu memberi masing-masing sebuah peluang yang sebanding lebih kecil, dan
kedua pengaruhnya meniadakan diri dengan tepat.
\end{solution}

\begin{solution}{exo:b3:molecular-evolution:2}
\omterm{def:b3:genomics:comparative}{Ortolog}: gen pada dua spesies yang berasal dari satu gen pada nenek
moyang bersamanya --- $\beta$-globin manusia dan $\beta$-globin mencit.
\omterm{def:b3:genomics:comparative}{Paralog}: gen di dalam satu genom yang berasal dari sebuah penggandaan ---
$\alpha$- dan $\beta$-globin manusia, atau $\beta$ dan $\gamma$. \omterm{def:b3:molecular-evolution:duplication}{Pseudogen}: sebuah
salinan yang membusuk yang tidak lagi menyandi sebuah protein --- yakni gen $\psi\beta$
di dalam gugusan $\beta$-globin manusia, lengkap dengan kodon hentinya.
\end{solution}

\begin{solution}{exo:b3:molecular-evolution:3}
$\omega$ membandingkan \omterm{thm:b3:molecular-evolution:neutral}{laju substitusi} yang mengubah asam amino dengan
\omterm{thm:b3:molecular-evolution:neutral}{laju substitusi} yang bisu, yang menjadi dasar netralnya. $0.02$:
\omterm{def:b3:molecular-evolution:dnds}{seleksi pemurni} yang kuat, dengan \qty{98}{\%} perubahan asam aminonya disingkirkan
--- aktin, \omterm{def:b3:chromatin-epigenetics:nucleosome}{histon} H3. $1.0$: tanpa kekangan --- sebuah \omterm{def:b3:molecular-evolution:duplication}{pseudogen}. $2.5$:
\omterm{def:b3:molecular-evolution:dnds}{seleksi positif}, dengan perubahan asam aminonya dipihak --- alur pengikat
peptida MHC, \omterm{def:b3:virology:virus}{selubung} HIV, lisozim hewan pemamah biak.
\end{solution}

\begin{solution}{exo:b3:molecular-evolution:4}
Protein yang dibandingkan di seluruh mamalia berubah pada sekitar satu substitusi tiap
tapak tiap miliar tahun; atas sebuah genom bertapak penyandi semiliar dan sebuah
angkatan beberapa tahun, itu berorde satu substitusi yang
terpancang tiap angkatan, atau satu tiap dua tahun. Ongkos seleksi
Haldane mengizinkan sekitar satu substitusi terpilih tiap tiga ratus
angkatan. Lajunya berbeda sebesar sebuah faktor seratus atau lebih, sehingga
nyaris semua perubahan yang terpancang pastilah netral.
\end{solution}

\begin{solution}{exo:b3:molecular-evolution:5}
Netralnya: $1/2N = 10^{-4}$. $s = 0.001$, $4Ns = 20$: $(1 - \mathrm{e}^{-0.002})/
(1 - \mathrm{e}^{-20}) = 0.002$. $s = 0.01$: $0.0198$. $s = -0.001$: $(1 -
\mathrm{e}^{0.002})/(1 - \mathrm{e}^{20}) = 0.002\,\mathrm{e}^{-20} =
4\times 10^{-12}$. Tidak satu pun netral efektif: itu akan menuntut $|4Ns|
\lesssim 1$, yakni $|s| \lesssim 5\times 10^{-5}$; sebuah koefisien
seleksi sepersepuluh persen sudah menentukan di dalam sebuah
populasi lima ribu.
\end{solution}

\begin{solution}{exo:b3:molecular-evolution:6}
Tanpa koreksi: $T = 0.12/(2\times 2\times 10^{-9}) = \qty{30}{Myr}$.
Jukes--Cantor: $d = -\tfrac{3}{4}\ln(1 - 0.16) = 0.131$, $T =
\qty{33}{Myr}$. Koreksinya mencapai sebuah faktor dua ketika $-\tfrac{3}{4}
\ln(1 - 4p/3) = 2p$, yakni pada $p \approx 0.6$: pada saat itu selisih mentahnya telah
kehilangan kebanyakan keterangannya.
\end{solution}

\begin{solution}{exo:b3:molecular-evolution:7}
$p_{N} = 7/700 = 0.010$, $p_{S} = 10/200 = 0.050$, $\omega = 0.2$:
yakni \omterm{def:b3:molecular-evolution:dnds}{seleksi pemurni} yang menyingkirkan empat perlima perubahan asam aminonya. Dengan
$35$ dan $10$: $p_{N} = 0.05$, $\omega = 1$: tanpa kekangan yang terdeteksi ---
yakni sebuah \omterm{def:b3:molecular-evolution:duplication}{pseudogen}, atau sebuah gen yang di dalamnya \omterm{def:b3:molecular-evolution:dnds}{seleksi positif} pada sebagian tapaknya
mengimbangi \omterm{def:b3:molecular-evolution:dnds}{seleksi pemurni} pada tapak yang lain.
\end{solution}

\begin{solution}{exo:b3:molecular-evolution:8}
$k = d/2T = 0.012/(1.3\times 10^{7}) = 9.2\times 10^{-10}$ tiap tapak tiap
tahun; $\mu = 25k = 2.3\times 10^{-8}$ tiap tapak tiap angkatan; $6\times
10^{9}\times 2.3\times 10^{-8} \approx 140$ mutasi baru tiap genom
diploid tiap angkatan. Pengurutan langsung orang tua dan anaknya menemukan
sekitar tujuh puluh: taksiran penyimpangannya dibesarkan polimorfisme
yang sudah hadir di dalam populasi leluhurnya, yang menambahkan beberapa ratus
ribu angkatan pada perpisahan yang tampak.
\end{solution}

\begin{solution}{exo:b3:molecular-evolution:9}
$\mathrm{e}^{-T/2N} = \mathrm{e}^{-1} = 0.37$; pecahan yang tak sesuainya $\tfrac{2}{3}
\times 0.37 = 0.25$. Bagi \qty{10}{\%}: $\mathrm{e}^{-T/2N} = 0.15$, $T =
2N\ln(1/0.15) = \num{190000}$ angkatan.
\end{solution}

\begin{solution}{exo:b3:molecular-evolution:10}
Ketika $s \to 0$: $1 - \mathrm{e}^{-2s} \approx 2s$ dan $1 - \mathrm{e}^{-4Ns}
\approx 4Ns$, nisbahnya $1/2N$. Bagi $4Ns \gg 1$ penyebutnya $1$ dan
pembilangnya $\approx 2s$. Bagi $s < 0$ dengan $4N|s| \gg 1$: pembilangnya $1 -
\mathrm{e}^{2|s|} \approx -2|s|$, penyebutnya $1 - \mathrm{e}^{4N|s|} \approx
-\mathrm{e}^{4N|s|}$, sehingga $P \approx 2|s|\,\mathrm{e}^{-4N|s|}$. Seratus kali lipat
di bawah netralnya pada $N = 10^{4}$: $2|s|\,\mathrm{e}^{-4\times 10^{4}|s|} =
5\times 10^{-7}$; $|s| = 1.6\times 10^{-4}$ memberi $3.2\times 10^{-4}\times
\mathrm{e}^{-6.4} = 5.3\times 10^{-7}$. Jadi $|s| \approx 1.6\times 10^{-4}$,
$4N|s| \approx 6.5$: sebuah kerugian seperseratus persen sudah
cukup, pada sepuluh ribu individu, untuk membuat pemancangannya seratus kali
lebih jarang daripada kebetulannya.
\end{solution}

\begin{solution}{exo:b3:molecular-evolution:11}
Dengan titik perpisahannya berjarak $a$ dari $O$ pada kedua jalannya, $a + m =
0.30$, $a + h = 0.24$, $m + h = 0.20$: $m - h = 0.06$, sehingga $m = 0.13$, $h
= 0.07$, dengan nisbahnya $1.9$. Sebuah jam tiap angkatan yang ketat akan memberi
nisbah cacah angkatannya, yakni $25/0.5 = 50$. Nilai $1.9$ yang teramati berarti
mencitnya telah menimbun hanya dua kali perubahan manusianya tiap tahun meskipun
angkatannya lima puluh kali lebih banyak: laju mutasi tiap angkatannya pastilah
sekitar $25$ kali lebih tinggi pada manusia --- yakni lebih banyak pembelahan sel nutfah tiap
angkatan --- sehingga tiap tahun kedua garis keturunannya berbeda hanya sebesar sebuah faktor
dua.
\end{solution}

\begin{solution}{exo:b3:molecular-evolution:12}
Substitusi yang menciptakan henti tiba pada $1000\times 2\times 10^{-9}\times
0.04 = 8\times 10^{-8}$ tiap tahun: yang pertama diharapkan sesudah sekitar
\qty{12}{Myr} (geseran kerangkanya kira-kira membagi duanya). Dalam jendela itu sebuah
mutasi yang menguntungkan, yakni satu dari $10^{3}$ sampai $10^{5}$ semua mutasinya, jauh
kurang mungkin tiba daripada sebuah kodon henti, dan bahkan ketika ia tiba ia
terpancang dengan peluang hanya $2s$. Jadi salinannya biasanya sudah mati sebelum
seleksi dapat menemukan sebuah kegunaan baginya --- paruh hidup gandaan yang teramati pada
vertebrata adalah beberapa juta tahun.
\end{solution}

\begin{solution}{pb:b3:molecular-evolution:1}
\textbf{1.} $k = 0.012/(2\times 6.5\times 10^{6}) = 9.2\times 10^{-10}$ tiap
tapak tiap tahun; $\mu = 25k = 2.3\times 10^{-8}$ tiap tapak tiap angkatan.
\textbf{2.} $6\times 10^{9}\times 2.3\times 10^{-8} \approx 140$ mutasi
baru tiap genom diploid; \qty{1.5}{\%} di antaranya, yakni sekitar $2$, di dalam
urutan penyandinya.
\textbf{3.} Tiap genom haploid, $3\times 10^{9}\times 2.3\times 10^{-8}
\approx 70$ substitusi terpancang tiap angkatan, terhadap $1/300$
milik Haldane: yakni dua puluh ribu kali lebih banyak daripada yang dapat digerakkan seleksi, sehingga
sebagian besarnya netral.
\textbf{4.} $2N\mu = 10^{5}\times 2.3\times 10^{-8} = 2.3\times 10^{-3}$
mutasi baru tiap tapak tiap angkatan di dalam populasinya; $1/2N =
10^{-5}$ di antaranya akan terpancang.
\textbf{5.} $4N = \num{200000}$ angkatan, yakni \qty{5}{Myr} --- nyaris
selama waktu sejak perpisahannya.
\textbf{6.} Penyimpangannya mencacah mutasi yang timbul sepanjang kedua
garis keturunan yang terpisah sejak nenek moyang bersamanya; masing-masing terpancang di dalam
garis keturunannya sendiri, dan laju jangka panjangnya $\mu$ berapa pun waktu yang dipakai
masing-masing, sebab mutasinya timbul mantap dan keterlambatannya hanya menunda, tanpa
mengubah, alirannya. (Satu-satunya koreksinya adalah bagi polimorfisme yang hadir pada
perpisahannya, yang membuat nenek moyang bersama dua alel lebih tua daripada
perpisahan spesiesnya.)
\textbf{7.} Netralnya $10^{-5}$; $s = 0.001$: $0.002$; $s = 0.01$: $0.0198$;
$s = 0.1$: $1 - \mathrm{e}^{-0.2} = 0.18$.
\textbf{8.} $s = 1/4N = 5\times 10^{-6}$: di bawah ini hanyutannya menenggelamkan
seleksinya dan mutasinya berlaku sebagai netral; di atasnya seleksinya
mengatur peluangnya.
\textbf{9.} $s = -0.001$: $P = 0.002\,\mathrm{e}^{-200}$, yakni nol bagi segala
maksud --- $10^{-85}$ dari netralnya. $s = -10^{-5}$, $4N|s| = 2$: $P =
2\times 10^{-5}/(\mathrm{e}^{2} - 1) = 3.1\times 10^{-6}$, yakni \qty{31}{\%} dari
netralnya: mutasi yang sedikit merugikan memang terpancang.
\textbf{10.} Substitusi yang menguntungkan tiap tapak tiap angkatan: $2N\times
10^{-5}\mu\times 2s = 10^{5}\times 10^{-5}\times 0.02\,\mu = 0.02\mu$;
netralnya: $\approx\mu$. Pecahan adaptifnya $0.02/1.02 \approx \qty{2}{\%}$.
Meskipun tiap mutasi yang menguntungkan dua ribu kali lebih mungkin terpancang
daripada mutasi netral, keduanya begitu jarang sehingga hanya satu substitusi dari
lima puluh yang bersifat adaptif: seleksinya bekerja dan nyaris tidak tampak.
\textbf{11.} $p_{N} = 46/920 = 0.050$, $p_{S} = 84/280 = 0.30$. Terkoreksi:
$d_{N} = -\tfrac{3}{4}\ln(1 - 0.0667) = 0.052$, $d_{S} = -\tfrac{3}{4}
\ln(1 - 0.40) = 0.38$. $\omega = 0.14$.
\textbf{12.} \omterm{def:b3:molecular-evolution:dnds}{Seleksi pemurni} yang menyingkirkan sekitar \qty{86}{\%}
perubahan asam aminonya --- yakni sebuah gen yang lazim. $k = d_{S}/2T = 0.38/(1.8\times
10^{8}) = 2.1\times 10^{-9}$ tiap tapak tiap tahun: yakni dua kali nilai
manusia--simpansenya, sebagaimana diharapkan bagi sebuah perbandingan yang menyertakan garis keturunan
hewan pengerat yang berdetak cepat; seorde besaran yang sama.
\textbf{13.} $d = -\ln(1 - 0.55) = 0.80$ tiap tapak; $T = 0.80/(2\times
10^{-9}) = \qty{400}{Myr}$.
\textbf{14.} Dua rantai yang berbeda membuat tetramer $\alpha_{2}\beta_{2}$,
yang pengikatan oksigennya yang kooperatif, \omterm{prop:b3:structural-biology:mechanism}{efek Bohr}, serta kendali alosterik
oleh 2,3-bisfosfogliseratnya bergantung pada antarmuka antara
subunit yang tak serupa; dan gugusan $\beta$-nya lalu dapat beraneka menjadi
rantai embrio, janin, dan dewasa dengan afinitas yang berbeda.
\textbf{15.} $2kT = 2\times 8\times 10^{-9}\times 4\times 10^{7} = 0.64$
substitusi tiap tapak; dengan kejenuhan di antara dua puluh asam aminonya,
selisih mentahnya sekitar $0.95(1 - \mathrm{e}^{-0.64/0.95}) \approx 0.47$.
Pada \qty{500}{Myr}, $d = 8$: tiap tapaknya telah berubah berkali-kali dan
selisih mentahnya duduk pada langit-langitnya \qty{95}{\%}, tanpa mengusung
keterangan tentang waktunya.
\textbf{16.} $10^{-11}\times 10^{9}\times 100\times 2 = 2$ substitusi
antara dua garis keturunan pada 100 tapak sepanjang semiliar tahun. Bagi sebuah
perpisahan \qty{10}{Myr} harapannya $0.02$: yakni nyaris selalu nol
selisih, tanpa daya pisah.
\textbf{17.} $1200\times 2\times 10^{-9}\times 0.04 = 9.6\times 10^{-8}$
tiap tahun: hentinya yang pertama diharapkan sesudah sekitar \qty{10}{Myr}.
\textbf{18.} Melipatduakan genom utuhnya melipatduakan tiap gennya sekaligus, sehingga
stoikiometri protein yang berinteraksinya terpelihara dan tidak ada ketakseimbangan
dosis yang menyeleksi melawan salinannya; sebuah gandaan tandem yang sendirian mengubah
dosis satu gennya lalu kerap merugikan, atau segera berlebihan dan
bebas membusuk.
\textbf{19.} $\mathrm{e}^{-\num{80000}/\num{100000}} = \mathrm{e}^{-0.8} = 0.45$.
\textbf{20.} $\tfrac{2}{3}\times 0.45 = 0.30$: sesuai dengan
\qty{30}{\%} yang teramati.
\textbf{21.} Separuhnya masing-masing, yakni \qty{15}{\%} semua pohonnya mengelompokkan manusia dengan
gorila dan \qty{15}{\%} simpanse dengan gorila. Sebuah kelebihan yang jelas pada salah satunya
akan menunjukkan aliran gen antara kedua spesies itu sesudah perpisahannya,
yang tidak dapat dihasilkan hanyutannya.
\textbf{22.} $N = \num{10000}$: $\mathrm{e}^{-4} = 0.018$, dengan ketaksesuaiannya
\qty{1.2}{\%}. Nilai \qty{30}{\%} yang teramati karena itu menuntut sebuah populasi
leluhur sekitar lima puluh ribu --- yakni beberapa kali ukuran efektif
manusia hari ini.
\textbf{23.} Perangkaiannya menganggap tiap gennya punya \omterm{prop:b3:molecular-evolution:ils}{pohon spesiesnya}; dengan
pemilahan garis keturunan yang tak lengkap, gennya punya banyak pohon, dan bagi cabang
dalam yang pendek, \omterm{prop:b3:molecular-evolution:ils}{pohon gen} yang terlazim dapat berbeda dari \omterm{prop:b3:molecular-evolution:ils}{pohon
spesiesnya}, sehingga lebih banyak data membuat jawaban yang salah lebih meyakinkan. Sebuah metode
koalesen menghitung, bagi tiap \omterm{prop:b3:molecular-evolution:ils}{pohon spesies} calonnya, harapan
sebaran \omterm{prop:b3:molecular-evolution:ils}{pohon gennya}, lalu memilih \omterm{prop:b3:molecular-evolution:ils}{pohon spesies} yang paling baik
menjelaskan kekerapan yang teramati.
\textbf{24.} Pemilahan garis keturunan yang tak lengkap terjadi di dalam populasi leluhur
bersama semua manusia modern, sehingga ia akan membuat tiap populasi
manusia sama berkerabatnya dengan Neanderthal. Sebuah kelebihan pada orang non-Afrika
berarti alel Neanderthalnya masuk sesudah nenek moyang orang non-Afrika
meninggalkan Afrika: yakni percampuran, sekitar lima puluh ribu tahun lalu.
\textbf{25.} $\mu \approx 2.3\times 10^{-8}$ tiap tapak tiap angkatan;
$\omega \approx 0.14$; sekitar \qty{30}{\%} \omterm{prop:b3:molecular-evolution:ils}{pohon gen} tidak sesuai dengan
\omterm{prop:b3:molecular-evolution:ils}{pohon spesiesnya}.
\end{solution}
